\bookchapter{第二十五章：二手手机处理与数字遗产}

\chapterintro{嘿，伙计们，经过前面二十多章的“修炼”，相信你已经从一个“数字小白”成长为一名相当从容的手机使用者了。但科技生活总有轮回——新欢终会变旧爱。当你准备让老伙计“退役”时，你知道该如何体面、安全地告别吗？}

\booksection{一、旧手机的“安乐死”：安全出售或转赠指南}

核心原则：数据清理不是“删除”，而是“覆盖”。许多人以为把照片拖进回收站、清空相册就万事大吉，结果手机转手后，新主人用数据恢复软件一扫，老照片、聊天记录、身份证照片全回来了——这种“社死现场”每年都在上演。

\begin{enumerate}
  \item \textbf{自我"大扫除"}：完成最后一次完整备份，退出所有账号（iCloud、华为/小米/三星账号、Google账号），关闭"查找我的手机"功能——否则下一位使用者会被激活锁拦在门外。
  \item \textbf{执行终极"格式化"}：现代手机的"恢复出厂设置"通常包含加密擦除——它销毁的是加密密钥本身，密钥一没，整盘数据瞬间变为乱码。
  \item \textbf{最后的"障眼法"（仅老设备需要）}：现代手机（iOS 8+、Android 6+）已默认全盘加密（FBE），恢复出厂设置会销毁加密密钥，数据理论上不可恢复，\textbf{无需此步骤}。仅对极老旧的非加密设备（2015年前），才需要在恢复出厂设置后用垃圾文件填满存储，再执行一次恢复出厂设置。
\end{enumerate}

\textbf{各品牌数据擦除路径速查表}：尽管原理相通，但每个品牌把“恢复出厂设置”藏的位置略有不同，关键差异如下。

\begin{center}
\begin{tabularx}{\linewidth}{|X|X|X|}
\hline
\textbf{品牌} & \textbf{擦除入口} & \textbf{关键差异与注意事项} \tabularnewline
\hline
华为 & 设置 > 系统和更新 > 重置 > 恢复出厂设置 & 需先退出华为账号并关闭"查找我的手机"，否则会触发激活锁；纯净模式下可能需要密码验证。 \tabularnewline
小米 / Redmi & 设置 > 我的设备 > 恢复出厂设置 & 底部有"格式化虚拟SD卡"和"清除全部内容与设置"选项，二者都要勾选；退出小米账号前先关闭"查找设备"。 \tabularnewline
OPPO / 一加 & 设置 > 其他设置 > 还原手机 > 彻底清除全部数据 & ColorOS 提供"彻底清除"模式，会多次覆盖存储；需输入锁屏密码确认。 \tabularnewline
vivo & 设置 > 系统管理 > 备份与重置 > 清除所有数据 & 勾选"格式化手机存储"才会清理SD卡区；需退出vivo账号，关闭"查找手机"。 \tabularnewline
三星 & 设置 > 通用管理 > 重置 > 恢复出厂设置 & 需先移除Samsung账号并关闭"查找我的手机"；部分机型有"自动覆盖"选项，建议开启。 \tabularnewline
iPhone & 设置 > 通用 > 传输或还原iPhone > 抹掉所有内容和设置 & 必须先退出Apple ID，否则下一位用户被iCloud激活锁拦住；iOS 15+支持"退出时保留iCloud数据"以备后查。 \tabularnewline
\hline
\end{tabularx}
\end{center}

\begin{tcbwarn}{最容易被忽略的三件事}
\begin{itemize}
  \item \textbf{未退出账号}：手机卖给下家后被远程锁定，自己也下不来，最后只能持身份证去官方售后解锁。
  \item \textbf{SD卡未取出}：恢复出厂设置不会格式化外置SD卡里的照片，记得手动取出或单独格式化。
  \item \textbf{U盘/OTG外设未断开}：擦除过程中如果连着外设，可能误清外设数据，也可能导致擦除中断。
\end{itemize}

\end{tcbwarn}

\booksection{二、数字遗产：你的社交灵魂如何安放？}

手机里的数据，早就不是“虚拟”的了。它承载着你的资金、社交关系、创作成果，乃至你的“数字身份”。\textbf{一个人去世后，这些资产该怎么处理？}这不再是一个抽象问题，而是越来越多家庭真实面临的难题——亲属不知道密码、平台不让继承、加密货币永远丢失。

\begin{tcbinfo}{法律依据：个人信息保护法第49条}
2021年施行的《中华人民共和国个人信息保护法》第49条明确规定：\textbf{自然人死亡的，其近亲属为了自身的合法、正当利益，可以对死者相关个人信息行使本章规定的查阅、复制、更正、删除等权利；死者生前另有安排的除外。}

这意味着：

\begin{itemize}
  \item \textbf{近亲属有法律依据}向平台申请处理死者的账号与个人信息。
  \item "合法、正当利益"通常包括办理遗产继承、关闭账号、删除敏感隐私等。
  \item 死者生前若有遗嘱或协议另行安排（如指定遗产联系人），从其约定。
  \item 平台不能以"用户隐私"为由一概拒绝配合，但具体流程各家不同。
\end{itemize}

\end{tcbinfo}

\booksubsection{微信与支付宝：账号不能继承，但资金可以}

这两个中国人最常用的App，规则各不相同：

\begin{itemize}
  \item \textbf{微信账号}：依据《腾讯微信软件许可及服务协议》，账号所有权归腾讯，用户仅享有使用权，\textbf{账号本身不能继承、不能转让}。但账号内的\textbf{零钱、零钱通资金}属于个人财产，近亲属可凭死亡证明、亲属关系证明、继承公证书向财付通申请提现或转移。
  \item \textbf{支付宝}：账号同样不能直接继承，但支付宝余额、余额宝、理财通等\textbf{资金属于可继承财产}。继承人可通过支付宝"客服热线"或线下提交死亡证明、继承公证书办理；蚂蚁财富类理财产品需走基金公司流程。
  \item \textbf{聊天记录}：法律上属个人隐私，平台不会主动提供给家属。如涉及遗产线索，可通过诉讼或公证申请调取。
\end{itemize}

\begin{tcbwarn}{千万别做的事}
\begin{itemize}
  \item \textbf{不要把死者账号密码直接交给家人"接着用"}——这违反平台协议，资金被冻结甚至触发风控时举证更难。
  \item \textbf{不要尝试用死者手机收验证码"提现"}——可能被认定为盗用，反而触犯法律。
  \item \textbf{不要拖延}：部分账号长期未登录会被自动注销，资金找回会更麻烦。
\end{itemize}

\end{tcbwarn}

\booksubsection{Apple“遗产联系人”：提前指定，事半而功倍}

iOS 15.2 起苹果推出“遗产联系人”功能，允许用户提前指定信任的联系人，去世后该联系人可凭\textbf{死亡证明 + 遗产联系人访问密钥}申请访问死者的iCloud数据（照片、备忘录、通讯录、备忘录等）。操作步骤：

\begin{enumerate}
  \item 在iPhone上进入\textbf{[设置] > [你的姓名] > [登录与安全] > [遗产联系人]}。
  \item 点击"添加遗产联系人"，使用面容ID确认身份。
  \item 选择联系人（从通讯录挑选），系统会生成一份\textbf{访问密钥}。
  \item 将访问密钥通过短信、邮件或打印纸质件交给对方，并提醒对方妥善保存。
  \item 可添加多个联系人作为备份，避免单点失效。
\end{enumerate}

被指定人无需Apple设备，凭密钥和死亡证明即可在苹果官网申请访问。整个流程通常在7天内完成。

\booksubsection{Google“不活跃账户管理器”：自动处理身后事}

Google提供“不活跃账户管理器（Inactive Account Manager）”，能在用户长期未登录时自动执行预设操作——这是目前最完善的数字遗产机制之一。

\begin{enumerate}
  \item 访问\textbf{myaccount.google.com/inactive}，登录Google账号。
  \item 设置超时期限：3 / 6 / 12 / 18个月未登录后触发。
  \item 添加受信任联系人（最多10人），并设置每位可访问哪些数据（Gmail、相册、云端硬盘等）。
  \item 可选：设置自动回复邮件，告知来信者你已无法处理。
  \item 最后选择：到期后是\textbf{分享数据给联系人}，还是\textbf{直接删除账户}。
  \item 设置手机号作为兜底验证，避免误触发。
\end{enumerate}

\booksubsection{国内社交账号：纪念化与注销流程}

国内主流社交平台对逝者账号的处理分两种思路：\textbf{“纪念化”（保留账号让亲友缅怀）}和\textbf{“注销”（彻底删除）}。家属可凭死亡证明、亲属关系证明、申请人身份证向平台客服申请。

\begin{itemize}
  \item \textbf{微博}：支持账号"纪念化"，由亲属提供死亡证明申请，账号会被锁定无法登录但内容保留，主页会显示"该账号已设置为纪念账号"。也可申请注销，需提供继承公证。
  \item \textbf{抖音}：通过"我 - 设置 - 反馈与帮助 - 账号与安全 - 账号注销"或人工客服，提交死亡证明、亲属关系证明、申请人身份证，审核通过后可注销或转为纪念状态。
  \item \textbf{B站}：客服渠道提交逝者证明及亲属证明，可申请账号注销或保留为纪念页；UP主的历史稿件处理需结合版权声明。
  \item \textbf{小红书}：通过官方客服申诉，提交死亡证明、亲属关系证明，可申请注销账号并删除内容。
\end{itemize}

\begin{tcbinfo}{关键提醒}
\begin{itemize}
  \item 各平台规则会更新，办理前先查阅官方客服最新指引。
  \item 所有证明文件建议先扫描备份再提交，原件妥善保管。
  \item 跨平台一次性办理时，可统一申请一份\textbf{继承公证书}，省去反复跑手续。
\end{itemize}

\end{tcbinfo}

\booksubsection{数字资产的继承难题}

除了社交账号，越来越多的人持有\textbf{真正有价值的数字资产}，其继承问题更加复杂：

\begin{itemize}
  \item \textbf{加密货币（比特币、以太坊等）}：基于"私钥即所有权"原则，\textbf{没有私钥就等于资产消失}。继承人无法通过法律文件向链上"申诉"。建议：将助记词/私钥分片存放在多个物理安全位置（如银行保险箱+家里保险柜），并在遗嘱中注明访问方式。
  \item \textbf{交易所账户}：和银行账户类似，可凭继承公证书向交易所（如币安、OKX）申请处置，但跨境交易所流程漫长。
  \item \textbf{域名}：属于可继承财产，继承人在域名注册商（如阿里云、GoDaddy）提交继承公证后可办理过户。
  \item \textbf{知识产权（原创文章、视频、专利、软件著作权）}：著作权中的财产权利可继承（人身权不继承），收益权可由继承人享有至作者死后50年。
  \item \textbf{云存储与网盘}：内容数据归属个人，可继承；但部分平台协议规定账号本身不可转让，需先下载再注销。
\end{itemize}

\textbf{加密货币“消失”的真实代价}\par

据Chainalysis报告，全球已有\textbf{超过300万枚比特币}因私钥丢失而永远无法找回，相当于现有总量的六分之一。如果你持有加密资产，请务必让至少一位你绝对信任的人知道\textbf{恢复路径}——否则这些财富将永远沉睡在区块链上。

\booksubsection{一份“数字遗嘱”模板}

传统的财产遗嘱管不了数字资产。建议你每年更新一份“数字遗嘱”附件，与正式遗嘱一并存放。以下模板供参考：

\textbf{《个人数字资产处置附件》（模板）}\par

\textbf{第一部分：账号清单}（仅记录账号名与所属平台，\textbf{不写密码}）

\begin{itemize}
  \item 社交类：微信（账号xxx）、微博、抖音、B站、小红书等
  \item 支付类：支付宝、微信支付、各银行卡App
  \item 云端类：iCloud、华为云、小米云、Google账号、微软账号
  \item 邮箱类：QQ邮箱、网易邮箱、Gmail等
  \item 其他：网盘、订阅服务、游戏账号等
\end{itemize}

\textbf{第二部分：密码托管方案}

\begin{itemize}
  \item 主密码存放方式：使用密码管理器（如Bitwarden、1Password），主密码写在纸质件中，封存于家中保险柜。
  \item 恢复密钥/助记词：分片存放，至少两处物理位置（如银行保管箱+家中保险柜）。
  \item 授权查询人：姓名、身份证号、与本人关系、联系方式。
\end{itemize}

\textbf{第三部分：处置意愿}

\begin{itemize}
  \item 各类账号希望"纪念化"还是"注销"。
  \item iCloud / Google 数据是否指定"遗产联系人"或"不活跃账户管理器"。
  \item 加密资产由谁继承，私钥获取路径说明。
  \item 版权类作品收益由谁享有。
\end{itemize}

\textbf{第四部分：联系人与平台}

\begin{itemize}
  \item 信任执行人姓名、联系方式、备用联系人。
  \item 常驻平台客服入口与所需证明材料清单。
  \item 本附件更新日期：\_\_\_\_年\_\_月\_\_日，下次复核日期：\_\_\_\_年\_\_月\_\_日。
\end{itemize}

\begin{tcbwarn}{关于"密码"的法律提示}
将密码直接写进遗嘱存在泄露风险，且遗嘱公示后所有继承人都能看到。\textbf{更稳妥的做法是：遗嘱里只写“账号清单 + 处置意愿 + 密码托管位置”，真正的密码通过密码管理器主密码单独传给信任执行人}。这样既符合法律要求，也降低了密码泄露风险。

\end{tcbwarn}

\booksection{三、环保的最后一程}

一部手机里含金、银、铜、稀土等数十种金属，随意丢弃既污染土壤，也浪费资源。正确的回收渠道有三类，各有侧重：

\begin{center}
\begin{tabularx}{\linewidth}{|X|X|X|X|}
\hline
\textbf{渠道} & \textbf{回收价} & \textbf{数据安全} & \textbf{适用场景} \tabularnewline
\hline
官方以旧换新（华为/苹果/小米等） & 中等，常叠加补贴券 & 最高，官方有规范流程 & 正打算换同品牌新机时最划算 \tabularnewline
第三方回收平台（爱回收、转转、回收宝等） & 略高，可比价 & 较高，有数据清除报告 & 想直接变现、不绑定品牌 \tabularnewline
公益捐赠（如乡村学校、爱心机构） & 无经济回报 & 需自行彻底擦除 & 设备仍可用、想做善事 \tabularnewline
发挥余热（自留改造） & 省下新设备钱 & 原机主本人继续用，无风险 & 设备性能尚可、家里有改造需求 \tabularnewline
\hline
\end{tabularx}
\end{center}

\begin{itemize}
  \item \textbf{官方回收/以旧换新}：最推荐的方式，流程规范，数据安全有保障。苹果、华为、小米等品牌官网都有"以旧换新"入口，评估后可抵扣新机款。
  \item \textbf{正规回收平台}：选择信誉良好的大型回收平台，下单后顺丰上门取件，检测后打款。务必选择\textbf{提供"数据清除报告"的}平台，并保留凭证。
  \item \textbf{公益捐赠}：可联系当地乡村学校、社区公益组织，旧手机可改装成教学平板或老人机使用。\textbf{捐赠前务必执行完整的擦除流程，并卸除所有SIM/SD卡}。
  \item \textbf{发挥余热}：改成车载导航仪、智能家居中控、监控摄像头、宝宝监视器、电子相框。一台旧手机可以再战好几年。
\end{itemize}

\begin{tcbinfo}{电子垃圾的真相}
一部手机的电路板里含有\textbf{金、银、钯、铜}等贵金属，每吨废旧电路板含金量是金矿石的数十倍。但若随意丢弃，其中的铅、汞、镉会渗入土壤与地下水，污染持续数十年。\textbf{正确回收，既保护环境，也能让稀缺金属循环再生。}

\end{tcbinfo}

\begin{tcbpractice}{实操：为你的旧手机执行一场"标准葬礼"}
\begin{enumerate}
  \item 最后一次完整备份：照片上传云端、重要文件复制到电脑、聊天记录导出存档。
  \item 退出所有账号：iCloud / 华为 / 小米 / Google 账号，关闭"查找我的手机"。
  \item 取出SIM卡与SD卡，断开所有OTG外设。
  \item 执行加密恢复出厂设置，必要时手动勾选"格式化存储"。
  \item （仅老设备）用垃圾文件填满存储，再恢复出厂设置一次。
  \item 选择它的下一站：官方以旧换新、正规回收平台、公益捐赠，或转赠家人。
  \item 同步更新你的"数字遗产清单"，把已淘汰的设备从清单中划掉。
\end{enumerate}

\end{tcbpractice}

\begin{tcbpractice}{本月任务清单}
\begin{enumerate}
  \item 为你目前的主力手机指定Apple"遗产联系人"或设置Google"不活跃账户管理器"。
  \item 用本节模板起草一份《个人数字资产处置附件》，与家人沟通存放位置。
  \item 检查家中是否有2台以上闲置旧手机，预约一次官方或正规平台回收。
\end{enumerate}

\end{tcbpractice}
